    \section{Выбранные методы визуально-инерциальной навигации}\label{sec:methods}

    \tikzset{
        blk/.style   = {draw, rounded corners=2pt, align=center, font=\scriptsize, minimum height=9mm, inner sep=3pt, text width=#1},
        blk/.default = 2.6cm,
        data/.style  = {draw, dashed, align=center, font=\scriptsize, minimum height=7mm, inner sep=3pt, text width=#1},
        data/.default = 2.2cm,
        arr/.style   = {-{Stealth[length=2mm]}, semithick},
        grp/.style   = {draw, dotted, rounded corners=3pt, inner sep=4pt}
    }

    В данной главе подробно рассматриваются три системы, выбранные в главе~\ref{sec:review} для экспериментального исследования: ORB-SLAM3, VINS-Fusion и OpenVINS.
    Для каждой системы описаны архитектура, обработка изображений, инициализация, способ оценки состояния, механизмы коррекции накопленной ошибки и практические особенности использования.
    Глава завершается сопоставлением систем и следствиями для постановки экспериментов.

    Ниже используются общие обозначения: $T_i = (R_i, p_i)$ -- поза аппарата в момент $i$, $v_i$ -- скорость, $b_{g,i}, b_{a,i}$ -- смещения гироскопа и акселерометра, $X_j$ -- положение $j$-й точки сцены.
    Ключевым кадром (keyframe) называется кадр, отобранный для участия в оптимизации или хранения в карте; остальные кадры используются только для отслеживания.

    \subsection{ORB-SLAM3}

    ORB-SLAM3~\cite{orbslam3} -- развитие систем ORB-SLAM~\cite{orbslam} и ORB-SLAM2, разработанное в Университете Сарагосы.
    Это полноценная SLAM-система: помимо оценки текущей позы она строит разреженную карту точек и ключевых кадров, распознаёт ранее посещённые места и корректирует по ним всю траекторию.
    В отличие от предшественников, ORB-SLAM3 тесно объединяет изображения с измерениями IMU и поддерживает несколько карт одновременно.

    \subsubsection{Архитектура}
    Система состоит из трёх параллельных потоков, работающих с общей структурой данных Atlas (рис.~\ref{fig:orbslam3}).

    \textbf{Atlas} хранит набор несвязанных карт: одну активную, в которой происходит локализация, и неактивные карты, построенные ранее.
    Каждая карта содержит ключевые кадры, точки и два графа: граф ковизибильности (рёбра между ключевыми кадрами с общими наблюдаемыми точками, вес ребра -- число общих точек) и основное дерево (spanning tree).
    Для распознавания мест используется база ключевых кадров на основе DBoW2~\cite{dbow2}, общая для всех карт.

    \textbf{Поток отслеживания} (tracking) обрабатывает каждый кадр и оценивает его позу относительно активной карты в реальном времени.
    Он же решает, нужно ли сделать текущий кадр ключевым.

    \textbf{Поток локального построения карты} (local mapping) добавляет ключевые кадры и новые точки в карту, удаляет избыточные элементы и уточняет локальный участок карты уравниванием связок.
    В визуально-инерциальном режиме этот же поток выполняет инициализацию IMU.

    \textbf{Поток замыкания петель и слияния карт} (loop closing and map merging) ищет совпадения нового ключевого кадра с ранее посещёнными местами в активной и неактивных картах и при обнаружении корректирует траекторию или объединяет карты.

    \begin{figure}[H]
        \centering
        \begin{tikzpicture}[node distance=5mm and 6mm]
            \node[data=1.8cm] (in) {Кадр\\+ измерения IMU};
            \node[blk, right=of in] (orb) {Извлечение ORB\\(пирамида масштабов)};
            \node[blk, right=of orb] (pred) {Прогноз позы\\(IMU или модель\\пост. скорости)};
            \node[blk=3.0cm, right=of pred] (lm) {Сопоставление\\с локальной картой,\\оптимизация позы};
            \node[grp, fit=(orb)(pred)(lm), label={[font=\scriptsize]above:Поток отслеживания}] (g1) {};

            \node[blk, below=12mm of lm] (map) {Новые точки,\\локальное BA,\\прореживание};
            \node[blk, left=of map] (imu) {Инициализация\\и уточнение IMU};
            \node[grp, fit=(map)(imu), label={[font=\scriptsize]below:Поток локального построения карты}] (g2) {};

            \node[blk=3.1cm, left=14mm of imu] (loop) {Распознавание места\\(DBoW2), замыкание\\петли / слияние карт};
            \node[grp, fit=(loop), label={[font=\scriptsize]above:Поток замыкания петель}] (g3) {};

            \node[data=3.1cm, below=12mm of loop] (atlas) {Atlas: активная\\и неактивные карты};

            \draw[arr] (in) -- (orb);
            \draw[arr] (orb) -- (pred);
            \draw[arr] (pred) -- (lm);
            \draw[arr] (lm) -- node[right, font=\scriptsize] {ключевой кадр} (map);
            \draw[arr] (map) -- (imu);
            \draw[arr] (imu) -- (loop);
            \draw[arr, dashed] (loop) -- (atlas);
            \draw[arr, dashed] (atlas.west) -| (in.south);
        \end{tikzpicture}
        \caption{Упрощённая схема ORB-SLAM3}
        \label{fig:orbslam3}
    \end{figure}

    \subsubsection{Обработка изображений и отслеживание}
    На каждом кадре строится пирамида изображений из восьми уровней с коэффициентом масштаба $1{,}2$, на каждом уровне детектором FAST выделяются углы, и для них вычисляются дескрипторы ORB~\cite{orb}.
    Чтобы точки распределялись равномерно, изображение разбивается на ячейки, и число точек в каждой ячейке ограничивается.
    Типичное число точек на кадр -- $1000$--$1500$.

    Поза нового кадра сначала прогнозируется: в визуальном режиме -- по модели постоянной скорости, в визуально-инерциальном -- интегрированием IMU от предыдущего кадра.
    Затем точки локальной карты (точки ключевых кадров, ковизибильных с текущим) проецируются в кадр и сопоставляются с его особыми точками по дескрипторам.
    По найденным соответствиям уточняется только поза текущего кадра (motion-only BA).
    В визуально-инерциальном режиме оптимизируются также скорость и смещения IMU, а в функцию стоимости, помимо ошибок перепроецирования, входит инерциальная невязка относительно предыдущего кадра.
    Если отслеживание теряется и восстановить его по базе ключевых кадров не удаётся, система начинает новую карту в Atlas, а не прекращает работу; позднее карты могут быть объединены.

    Ключевые кадры добавляются часто, а избыточные затем удаляются потоком локального построения карты: ключевой кадр считается избыточным, если $90\%$ его точек наблюдаются как минимум тремя другими ключевыми кадрами.
    Такая стратегия делает систему устойчивой к резким манёврам, при которых необходимо быстро расширять карту.

    \subsubsection{Инициализация}
    В монокулярном режиме карта инициализируется по двум кадрам с достаточным параллаксом.
    Параллельно оцениваются две модели -- гомография, соответствующая плоской сцене или чистому повороту, и фундаментальная матрица для произвольной сцены, -- и выбирается та, которая лучше объясняет соответствия.
    Для съёмки с БПЛА этот механизм существенен: при направленной вниз камере сцена близка к плоскости, и инициализация обычно происходит по гомографии.

    Инициализация IMU выполняется по методу~\cite{orbslam3_init} в три этапа.
    Сначала около двух секунд работает чисто визуальная система, строящая карту с точностью до масштаба.
    Затем решается задача оценки апостериорного максимума только по инерциальным измерениям при фиксированной визуальной траектории: находятся масштаб $s$, направление гравитации $R_{wg}$, скорости $v_i$ и смещения $b_g, b_a$.
    Наконец, все переменные уточняются совместным визуально-инерциальным уравниванием связок.
    Через $5$ и $15$ секунд после инициализации инерциальные параметры уточняются повторно.
    Ключевая особенность подхода -- учёт неопределённости IMU в виде априорной информации, что позволяет получить приемлемую оценку масштаба по малому числу ключевых кадров.
    Однако при равномерном прямолинейном полёте ускорения малы, масштаб и смещение акселерометра плохо различимы, и инициализация может затянуться или оказаться неточной.

    \subsubsection{Оптимизация}
    Локальное уравнивание связок выполняется в потоке локального построения карты после добавления каждого ключевого кадра.
    В визуальном режиме оптимизируются текущий ключевой кадр, все ковизибильные с ним ключевые кадры и наблюдаемые ими точки; ключевые кадры, видящие эти точки, но не входящие в окно, фиксируются.
    В визуально-инерциальном режиме окно образуют последние ключевые кадры, связанные цепочкой инерциальных ограничений, а оптимизируемыми переменными являются
    \begin{equation}\label{eq:orb_vi_state}
        \mathcal{S}_i = \left\{ T_i,\ v_i,\ b_{g,i},\ b_{a,i} \right\}, \qquad
        \min_{\{\mathcal{S}_k\},\, \{X_j\}} \left( \sum_{k} \left\| r_{\mathcal{I}_{k-1,k}} \right\|^2_{\Sigma_{\mathcal{I}}^{-1}} + \sum_{k}\sum_{j} \rho\left( \left\| r_{kj} \right\|^2_{\Sigma_{kj}^{-1}} \right) \right),
    \end{equation}
    где $r_{\mathcal{I}_{k-1,k}}$ -- невязка преинтегрированных измерений IMU~\cite{preintegration} между соседними ключевыми кадрами, $r_{kj}$ -- ошибка перепроецирования точки $X_j$ в кадр $k$, $\rho$ -- функция Хьюбера.
    Для решения используется библиотека g2o~\cite{g2o}.
    Камера может быть описана моделью pinhole или моделью Каннала--Брандта~\cite{kannala_brandt} для широкоугольных объективов; алгоритмы системы не зависят от конкретной модели.

    \subsubsection{Замыкание петель и слияние карт}
    Для каждого нового ключевого кадра по базе DBoW2 подбираются кандидаты на совпадение.
    В отличие от ORB-SLAM2, где требовалось подтверждение совпадения тремя последовательными ключевыми кадрами, ORB-SLAM3 сначала проверяет геометрическую согласованность кандидата с несколькими ковизибильными ключевыми кадрами, что ускоряет обнаружение петель и повышает их полноту.
    Найденное совпадение задаёт преобразование между текущей и найденной позой: $Sim(3)$ в монокулярном режиме, где масштаб ненаблюдаем, и $SE(3)$ при наличии стереопары или IMU.

    Если совпадение найдено в активной карте, выполняется замыкание петли: ошибка распределяется по траектории оптимизацией графа поз (essential graph), после чего в отдельном потоке запускается полное уравнивание связок.
    Если совпадение найдено в неактивной карте, карты объединяются: точки и ключевые кадры совмещаются, после чего уточняется «окно сварки» (welding window) вокруг места соединения.

    Для поисковых миссий БПЛА это свойство особенно ценно: при повторном пролёте над участком накопленная ошибка может быть частично устранена.
    Однако после замыкания петли положение меняется скачком, а распознавание мест по ORB-дескрипторам ненадёжно при съёмке одного участка с другой высоты или под другим курсом.

    \subsubsection{Практические особенности}
    Система написана на C++ и зависит от OpenCV, Eigen, Pangolin (визуализация), а также включает модифицированные копии DBoW2 и g2o.
    Распространяется по лицензии GPLv3.
    Репозиторий содержит примеры запуска на наборах EuRoC и TUM-VI без ROS и узлы для ROS~1; для ROS~2 существуют только сторонние обёртки.
    Результаты сохраняются в текстовые файлы форматов TUM и EuRoC: полная траектория -- для стерео- и инерциальных режимов, траектория ключевых кадров -- для монокулярного.
    Параметры камеры, IMU и алгоритма задаются YAML-файлом; калибровка камеры и IMU выполняется внешними средствами.

    \subsection{VINS-Mono и VINS-Fusion}

    VINS-Mono~\cite{vins_mono} -- система монокулярной визуально-инерциальной одометрии, разработанная в Гонконгском университете науки и технологии для БПЛА и мобильных устройств.
    VINS-Fusion~\cite{vins_fusion} обобщает её на стереокамеру, стереокамеру с IMU и добавляет модуль глобального объединения с GNSS, сохраняя ту же структуру оценивателя.
    Далее обе системы описываются вместе, так как в монокулярно-инерциальном режиме VINS-Fusion повторяет VINS-Mono.

    \subsubsection{Архитектура}
    Система построена как набор узлов ROS (рис.~\ref{fig:vins}):
    \begin{enumerate}
        \item \textbf{предобработка измерений} -- отслеживание особых точек и преинтегрирование IMU между кадрами;
        \item \textbf{инициализация} -- восстановление масштаба, гравитации, скоростей и смещения гироскопа;
        \item \textbf{локальная VIO} -- нелинейная оптимизация в скользящем окне с маргинализацией;
        \item \textbf{глобальная оптимизация} -- распознавание мест, релокализация и оптимизация графа поз;
        \item \textbf{глобальное объединение} (только VINS-Fusion) -- совмещение локальной траектории с измерениями GNSS.
    \end{enumerate}

    \begin{figure}[H]
        \centering
        \begin{tikzpicture}[node distance=5mm and 6mm]
            \node[data=1.6cm] (cam) {Кадр};
            \node[data=1.6cm, below=12mm of cam] (imu) {IMU};
            \node[blk=2.9cm, right=of cam] (ft) {Ши--Томаси + KLT,\\RANSAC по\\фундам. матрице};
            \node[blk=2.9cm, right=of imu] (pre) {Преинтегрирование\\IMU};
            \node[blk=2.6cm, right=8mm of ft, yshift=-10mm] (init) {Инициализация:\\SfM + визуально-\\инерциальное\\выравнивание};
            \node[blk=2.8cm, right=of init] (win) {Оптимизация\\в скользящем окне\\(Ceres),\\маргинализация};
            \node[blk=2.6cm, below=10mm of win] (pg) {DBoW2 + BRIEF,\\релокализация,\\граф поз 4 DoF};
            \node[blk=2.6cm, left=of pg] (gps) {Глобальное\\объединение\\с GNSS};
            \node[data=1.8cm, right=of win] (out) {Поза\\с частотой\\IMU};

            \draw[arr] (cam) -- (ft);
            \draw[arr] (imu) -- (pre);
            \draw[arr] (ft) -- (init);
            \draw[arr] (pre) -- (init);
            \draw[arr] (init) -- (win);
            \draw[arr] (win) -- (out);
            \draw[arr] (win) -- node[right, font=\scriptsize] {ключевые кадры} (pg);
            \draw[arr] (pg) -- (gps);
        \end{tikzpicture}
        \caption{Упрощённая схема VINS-Mono / VINS-Fusion}
        \label{fig:vins}
    \end{figure}

    \subsubsection{Обработка изображений}
    На каждом кадре отслеживаются точки предыдущего кадра пирамидальным методом Лукаса--Канаде~\cite{lucas_kanade}, а при нехватке точек добавляются новые углы Ши--Томаси~\cite{shi_tomasi} с ограничением на минимальное расстояние между ними (типично $100$--$200$ точек на кадр).
    Выбросы отсеиваются схемой RANSAC с оценкой фундаментальной матрицы.
    Дескрипторы на этом этапе не вычисляются, что делает предобработку значительно дешевле, чем в ORB-SLAM3, но лишает систему возможности восстанавливать соответствия после потери отслеживания.

    Кадр становится ключевым, если средний параллакс отслеживаемых точек относительно предыдущего ключевого кадра превышает порог или если число отслеживаемых точек опускается ниже порога.
    При вычислении параллакса используется компенсация поворота по гироскопу, иначе поворот на месте ошибочно порождал бы ключевые кадры.

    \subsubsection{Инициализация}
    Инициализация разделена на визуальную и инерциальную части.
    Сначала по кадрам окна выполняется восстановление структуры по движению (SfM) с точностью до масштаба: выбирается пара кадров с достаточным параллаксом, относительная поза находится по пяти точкам~\cite{five_point}, точки триангулируются, позы остальных кадров находятся решением PnP, и всё уточняется уравниванием связок.
    Затем визуальная траектория выравнивается с преинтегрированными измерениями IMU: последовательно оцениваются смещение гироскопа (по согласованию поворотов), скорости, вектор гравитации и масштаб (линейная задача наименьших квадратов), после чего гравитация уточняется с учётом её известной величины.
    Смещение акселерометра на этом этапе не оценивается и принимается нулевым -- оно уточняется уже в оптимизации.

    \subsubsection{Оптимизация в скользящем окне}
    Состояние окна включает $n+1$ ключевых кадров (по умолчанию $n = 10$), обратные глубины точек и параметры калибровки:
    \begin{equation}\label{eq:vins_state}
        \mathcal{X} = \left[ x_0, \dots, x_n,\ x_c^b,\ \lambda_1, \dots, \lambda_m \right], \qquad
        x_k = \left[ p_k,\ v_k,\ q_k,\ b_{a,k},\ b_{g,k} \right],
    \end{equation}
    где $x_c^b$ -- положение камеры относительно IMU, $\lambda_l$ -- обратная глубина точки $l$ в кадре её первого наблюдения.
    Параметризация обратной глубиной задаёт каждую точку одним числом вместо трёх и хорошо ведёт себя для далёких точек.
    Решается задача
    \begin{equation}\label{eq:vins_cost}
        \min_{\mathcal{X}} \left\{ \left\| r_p - H_p \mathcal{X} \right\|^2 + \sum_{k} \left\| r_{\mathcal{B}}\!\left( \hat{z}_{b_{k+1}}^{b_k}, \mathcal{X} \right) \right\|^2_{P_{b_{k+1}}^{b_k}} + \sum_{(l,j)} \rho\left( \left\| r_{\mathcal{C}}\!\left( \hat{z}_l^{c_j}, \mathcal{X} \right) \right\|^2_{P_l^{c_j}} \right) \right\},
    \end{equation}
    где первое слагаемое -- априорная информация, оставшаяся после маргинализации, второе -- невязки преинтегрированных измерений IMU между соседними кадрами окна, третье -- визуальные невязки, вычисляемые на единичной сфере, что позволяет использовать широкоугольные камеры.
    Задача решается библиотекой Ceres Solver~\cite{ceres}.

    Маргинализация поддерживает размер окна постоянным.
    Если второй по новизне кадр является ключевым, из окна удаляется самый старый кадр, и его измерения переводятся в априорное слагаемое через дополнение Шура.
    В противном случае удаляются визуальные измерения второго по новизне кадра, а его инерциальные измерения сохраняются: так окно содержит разнесённые ключевые кадры, но не теряет информацию IMU.

    \textbf{Онлайн-калибровка.}
    Положение камеры относительно IMU $x_c^b$ может уточняться в процессе работы.
    Кроме того, в VINS-Fusion оценивается временной сдвиг $t_d$ между метками времени камеры и IMU~\cite{vins_td}: каждая точка сдвигается по изображению на $t_d V_l$, где $V_l$ -- её скорость на изображении, и $t_d$ входит в оптимизацию как обычная переменная.
    Это особенно полезно при использовании недорогих несинхронизированных камер.

    \subsubsection{Релокализация и глобальная оптимизация}
    Для распознавания мест на ключевых кадрах дополнительно выделяется около $500$ углов FAST с дескрипторами BRIEF, которые используются только для поиска по базе DBoW2.
    Кандидаты на совпадение проверяются в два этапа: RANSAC по фундаментальной матрице для соответствий «2D--2D» и RANSAC в задаче PnP для соответствий «3D--2D».
    Подтверждённое совпадение добавляется в оптимизацию окна как дополнительное ограничение (релокализация), что сразу уменьшает дрейф текущей оценки.

    Затем выполняется глобальная оптимизация графа поз.
    Поскольку при использовании IMU крен и тангаж наблюдаемы, граф оптимизируется только по четырём степеням свободы -- положению и рысканию~\cite{vins_mono}.
    Это уменьшает размер задачи и не допускает искажения вертикали ошибочными петлями.
    Граф поз может сохраняться и загружаться повторно, что позволяет переиспользовать карту.

    \subsubsection{Объединение с GNSS}
    Модуль глобального объединения VINS-Fusion рассматривает результат VIO как локальные относительные измерения, а GNSS -- как глобальные абсолютные и оптимизирует граф поз, связывающий их.
    Для данной работы этот модуль представляет интерес как готовый механизм для сценария частичной деградации: при плохом GNSS его вес снижается, а при отсутствии траектория продолжается по VIO.
    Однако основная задача работы -- передача визуальной оценки в EKF автопилота, который сам выполняет объединение с GNSS, поэтому модуль используется только для сравнения.

    \subsubsection{Практические особенности}
    Система написана на C++, зависит от ROS~1 (Kinetic/Melodic, с доработками -- Noetic), OpenCV, Eigen и Ceres Solver, распространяется по лицензии GPLv3.
    Режим работы и параметры задаются YAML-файлом; камеры описываются через библиотеку camodocal, поддерживающую модели pinhole, MEI и Каннала--Брандта.
    Входные данные подаются через топики ROS, поэтому записанные данные необходимо преобразовать в формат rosbag.
    Траектория VIO и траектория после коррекции петлями публикуются раздельно и записываются в файлы, что удобно для оценки вклада замыкания петель.
    Известные недостатки -- чувствительность к качеству начальной калибровки и к параметрам шумов IMU, а также отсутствие официальной поддержки ROS~2.

    \subsection{OpenVINS}

    OpenVINS~\cite{openvins} -- открытая платформа для исследований визуально-инерциальной навигации, разработанная в Университете Делавэра.
    В её основе лежит фильтр MSCKF~\cite{msckf}, дополненный хранением долгоживущих точек в состоянии и онлайн-калибровкой.
    В отличие от ORB-SLAM3 и VINS-Fusion, задуманных как готовые системы, OpenVINS спроектирован как модульная библиотека с подробной документацией, симулятором и инструментами оценки.

    \subsubsection{Архитектура}
    Проект состоит из нескольких пакетов:
    \begin{itemize}
        \item \texttt{ov\_core} -- отслеживание особых точек, модели камер, математические утилиты;
        \item \texttt{ov\_init} -- статическая и динамическая инициализация;
        \item \texttt{ov\_msckf} -- оцениватель на основе MSCKF;
        \item \texttt{ov\_eval} -- инструменты оценки траекторий (совмещение, ATE, RPE, проверка консистентности по NEES);
        \item \texttt{ov\_data} -- конфигурации и эталонные траектории для открытых наборов данных.
    \end{itemize}
    В состав проекта также входит симулятор, порождающий синтетические измерения камеры и IMU вдоль заданной B-сплайновой траектории.
    Цикл работы фильтра показан на рис.~\ref{fig:openvins}.

    \begin{figure}[H]
        \centering
        \begin{tikzpicture}[node distance=5mm and 6mm]
            \node[data=1.6cm] (imu) {IMU};
            \node[data=1.6cm, below=10mm of imu] (cam) {Кадр};
            \node[blk=2.5cm, right=of imu] (prop) {Прогноз\\состояния\\и ковариации};
            \node[blk=2.5cm, right=of cam] (trk) {FAST + KLT\\или ORB,\\RANSAC};
            \node[blk=2.4cm, right=8mm of prop, yshift=-8mm] (clone) {Клонирование\\позы (stochastic\\cloning)};
            \node[blk=2.7cm, right=of clone, yshift=8mm] (msckf) {MSCKF-обновление:\\триангуляция,\\проекция на\\нуль-пространство};
            \node[blk=2.7cm, below=4mm of msckf] (slam) {SLAM-обновление\\долгоживущих точек};
            \node[blk=2.4cm, right=of msckf, yshift=-7mm] (marg) {Маргинализация\\старейшего\\клона};

            \draw[arr] (imu) -- (prop);
            \draw[arr] (cam) -- (trk);
            \draw[arr] (prop) -- (clone);
            \draw[arr] (trk) -- (clone);
            \draw[arr] (clone) -- (msckf);
            \draw[arr] (clone) -- (slam);
            \draw[arr] (msckf) -- (marg);
            \draw[arr] (slam) -- (marg);
            \draw[arr, dashed] (marg.north) |- ([yshift=5mm]msckf.north) -| node[pos=0.25, above, font=\scriptsize] {следующий кадр} (prop.north);
        \end{tikzpicture}
        \caption{Цикл обработки кадра в OpenVINS}
        \label{fig:openvins}
    \end{figure}

    \subsubsection{Вектор состояния}
    Состояние фильтра в момент $k$ имеет вид
    \begin{equation}\label{eq:ov_state}
        \mathbf{x}_k = \left[ \mathbf{x}_I^\top,\ \mathbf{x}_C^\top,\ \mathbf{x}_F^\top,\ \mathbf{x}_W^\top,\ t_d \right]^\top, \qquad
        \mathbf{x}_I = \left[ q_k^\top,\ p_k^\top,\ v_k^\top,\ b_{g,k}^\top,\ b_{a,k}^\top \right]^\top,
    \end{equation}
    где $\mathbf{x}_I$ -- состояние IMU, $\mathbf{x}_C$ -- скользящее окно клонов (копий поз IMU в моменты предыдущих кадров), $\mathbf{x}_F$ -- долгоживущие SLAM-точки, $\mathbf{x}_W$ -- параметры калибровки (внутренние параметры камеры и её положение относительно IMU), $t_d$ -- временной сдвиг.
    Какие из калибровочных параметров оцениваются, задаётся в конфигурации.

    \subsubsection{Прогноз и обновление}
    Между кадрами состояние IMU прогнозируется интегрированием измерений гироскопа и акселерометра, а ковариация -- по линеаризованной модели ошибок.
    При поступлении кадра текущая поза IMU копируется в окно клонов (stochastic cloning) с сохранением взаимных ковариаций.

    Особые точки отслеживаются либо методом Лукаса--Канаде по углам FAST, либо сопоставлением дескрипторов ORB; изображение разбивается на сетку для равномерного распределения точек.
    Точки используются для обновления двумя способами.

    \textbf{MSCKF-обновление.}
    Точка, отслеживание которой закончилось или которая наблюдалась во всех клонах окна, триангулируется по позам клонов.
    Её невязки линеаризуются по состоянию $\mathbf{x}$ и положению точки $X_j$:
    \begin{equation}\label{eq:msckf_lin}
        r_j = H_{x,j}\, \tilde{\mathbf{x}} + H_{f,j}\, \tilde{X}_j + n_j .
    \end{equation}
    Положение точки не входит в состояние, поэтому зависимость от него исключается умножением на матрицу $N_j^\top$, столбцы которой образуют базис левого нуль-пространства $H_{f,j}$:
    \begin{equation}\label{eq:msckf_null}
        N_j^\top r_j = N_j^\top H_{x,j}\, \tilde{\mathbf{x}} + N_j^\top n_j .
    \end{equation}
    Полученное ограничение связывает только позы клонов и используется в стандартном обновлении EKF.
    Благодаря этому вычисления растут линейно с числом точек, а не кубически, как при хранении всех точек в состоянии.

    \textbf{SLAM-обновление.}
    Точки, отслеживаемые дольше окна, добавляются в состояние как SLAM-точки и обновляются обычным образом при каждом новом наблюдении.
    Они дают ограничения на длительных интервалах, что снижает дрейф при движении над одним участком, например при зависании.

    После обновления самый старый клон удаляется из состояния вместе с соответствующими строками и столбцами ковариации.
    Отдельно поддерживается обновление нулевой скоростью (zero-velocity update) при обнаружении неподвижности, полезное перед взлётом и при посадке.

    \subsubsection{Консистентность}
    Особое внимание в OpenVINS уделено консистентности фильтра.
    При стандартной линеаризации EKF якобианы вычисляются в последних оценках, и система получает ложную информацию о ненаблюдаемых направлениях (положение и рыскание); фильтр становится излишне уверенным в своей оценке~\cite{vins_consistency}.
    OpenVINS использует метод якобианов первой оценки (First-Estimate Jacobians)~\cite{fej}: для каждой переменной якобианы вычисляются в её первой оценке, что сохраняет правильную размерность ненаблюдаемого подпространства.
    Для данной работы это важно: ковариация, которую выдаёт фильтр, используется EKF автопилота как вес измерения, и заниженная ковариация может ухудшить общее решение.

    \subsubsection{Инициализация}
    Поддерживаются два способа инициализации.
    Статическая требует, чтобы аппарат некоторое время был неподвижен: по средним показаниям акселерометра находится направление гравитации, по гироскопу -- его смещение, а начало движения обнаруживается по скачку ускорения.
    Для БПЛА это естественный вариант -- инициализация на земле перед взлётом.
    Динамическая инициализация решает линейную задачу по нескольким кадрам в движении и позволяет запускать фильтр в полёте, например после сбоя.

    \subsubsection{Практические особенности}
    Система написана на C++, зависит от OpenCV, Eigen, Boost и Ceres Solver (для динамической инициализации), поддерживает ROS~1 и ROS~2 и распространяется по лицензии GPLv3.
    Поддерживаются монокулярная, стерео- и многокамерная конфигурации, модели камер pinhole и equidistant (fisheye).
    Замыкание петель в основной системе отсутствует; оно возможно только через сторонние надстройки, использующие модуль петель VINS-Fusion.
    Основное преимущество для исследования -- встроенные инструменты оценки и симулятор, позволяющий проверить корректность настройки на синтетических данных до работы с реальными.

    \subsection{Сопоставление систем}

    Основные характеристики систем сведены в таблицу~\ref{table:selected_methods}.

    \begin{table}[H]
        \caption{Сопоставление ORB-SLAM3, VINS-Fusion и OpenVINS}
        \label{table:selected_methods}
        \scriptsize
        \centering
        \renewcommand{\arraystretch}{1.25}
        \renewcommand\tabularxcolumn[1]{m{#1}}
        \begin{tabularx}{\textwidth}{|>{\centering\arraybackslash}m{0.17\textwidth}|>{\centering\arraybackslash}X|>{\centering\arraybackslash}X|>{\centering\arraybackslash}X|}
            \hline
            \textbf{Характеристика} & \textbf{ORB-SLAM3} & \textbf{VINS-Fusion} & \textbf{OpenVINS} \\
            \hline
            Класс & SLAM, оптимизация & VIO, оптимизация в окне & VIO, фильтр MSCKF \\
            \hline
            Конфигурации датчиков & моно, стерео, RGB-D; каждая с IMU и без & моно + IMU, стерео, стерео + IMU, + GNSS & моно, стерео, многокамерная; только с IMU \\
            \hline
            Особые точки & FAST + ORB, сопоставление дескрипторов & Ши--Томаси + KLT & FAST + KLT или ORB \\
            \hline
            Инициализация IMU & визуальная карта + инерциальная MAP-оценка & SfM + визуально-инерциальное выравнивание & статическая или динамическая \\
            \hline
            Замыкание петель & да, $Sim(3)/SE(3)$, слияние карт & да, граф поз 4 DoF & нет \\
            \hline
            Онлайн-калибровка & нет & положение камеры, временной сдвиг & внутренние параметры, положение камеры, временной сдвиг \\
            \hline
            Потеря отслеживания & новая карта в Atlas & повторная инициализация & повторная инициализация \\
            \hline
            Ковариация на выходе & нет & нет (только граф) & да, консистентная (FEJ) \\
            \hline
            Решатель & g2o & Ceres Solver & собственная реализация EKF \\
            \hline
            ROS & ROS~1 (ROS~2 -- сторонние обёртки) & ROS~1 & ROS~1 и ROS~2 \\
            \hline
            Лицензия & GPLv3 & GPLv3 & GPLv3 \\
            \hline
        \end{tabularx}
    \end{table}

    Из сопоставления следуют выводы для постановки экспериментов.
    Во-первых, системы различаются по входным данным: ORB-SLAM3 работает с файлами и ROS~1, VINS-Fusion -- только с ROS~1, OpenVINS -- с ROS~1 и ROS~2.
    Поэтому для воспроизводимого сравнения необходим общий конвейер подготовки данных, приводящий запись или симуляцию к формату EuRoC и rosbag, и единый формат выходных траекторий (TUM) для оценки.
    Во-вторых, только ORB-SLAM3 и VINS-Fusion корректируют ошибку замыканием петель, поэтому в экспериментах необходимо отдельно оценивать траекторию до и после коррекции, а маршруты должны включать как прямолинейные участки, так и повторные пролёты.
    В-третьих, для передачи в автопилот важны выходная ковариация и поведение при потере отслеживания: OpenVINS выдаёт согласованную ковариацию, а для двух других систем её придётся задавать эмпирически по результатам экспериментов.
    Наконец, все три системы чувствительны к инерциальной инициализации, а условия полёта с малыми ускорениями неблагоприятны для неё; это требует проверки в сценариях с различной динамикой.

    \subsection{Выводы по главе}

    Выбранные системы представляют три разных подхода к визуально-инерциальной навигации.
    ORB-SLAM3 -- наиболее универсальная SLAM-система с сопоставлением дескрипторов, замыканием петель и множеством карт, устойчивая к потере отслеживания, но наиболее требовательная к текстуре сцены и не выдающая ковариацию.
    VINS-Fusion -- оптимизационная VIO с дешёвым отслеживанием потоком, онлайн-калибровкой временного сдвига и глобальной коррекцией по четырём степеням свободы, изначально разработанная для БПЛА.
    OpenVINS -- фильтрационная VIO с линейной по числу точек сложностью, консистентной оценкой ковариации и развитыми средствами оценки, но без замыкания петель.

    Различия в архитектуре определяют гипотезы, проверяемые в экспериментах: ожидается, что ORB-SLAM3 и VINS-Fusion дадут наименьшую ошибку на маршрутах с повторными пролётами, OpenVINS -- наименьшие вычислительные затраты и наиболее предсказуемый рост ошибки, а на однородных поверхностях преимущество получат методы с отслеживанием потоком (VINS-Fusion, OpenVINS) над методом на ORB-дескрипторах.
    \clearpage
